% Full title: Feasibility Study of Direct DtN Maps for Periodic Half-Guides
\documentclass[11pt]{article}
\usepackage[margin=1in]{geometry}
\usepackage{amsmath,amssymb}
\usepackage{hyperref}
\usepackage{booktabs}

\title{Feasibility Study: Direct DtN Maps for Periodic Half-Guides in the TEP Formulation}
\author{TEP DtN Attempt}
\date{Draft status: July 13, 2026}

\begin{document}
\maketitle

\section*{Executive Summary}

This is a draft feasibility report created at the beginning of the isolated
\texttt{attempt/} study. The current evidence supports the sign conversion for
the homogeneous lead, but does not yet validate obstacle leads or the full
Muller-DtN eigenvalue formulation.

\section{Current Formulation And Motivation}

The current project uses outgoing Bloch trace subspaces at the two center-cell
ports. The proposed alternative is to replace these subspaces by half-guide
Dirichlet-to-Neumann maps
\[
  N^\pm=\Lambda^\pm(k,\beta)D^\pm,
  \qquad
  N^+=+\partial_xu|_{\Gamma^+},\quad
  N^-=-\partial_xu|_{\Gamma^-}.
\]

\section{Review Of Joly And Coatleven}

Joly, Li, and Fliss construct exact boundary conditions for periodic
waveguides using local cell DtN blocks and a stationary Riccati equation
\cite{JolyLiFliss2006}. Coatleven reuses this construction for line defects
after a Floquet-Bloch transform along the defect direction \cite{Coatleven2012}.

\section{Homogeneous Lead Check}

For one Rayleigh channel in a homogeneous cell of length \(L\),
\[
  T_{00}=T_{11}=\gamma\cot(\gamma L),\qquad
  T_{01}=T_{10}=-\gamma\csc(\gamma L).
\]
With \(R=\exp(i\gamma L)\), the Joly Riccati residual vanishes:
\[
  T_{10}R^2+(T_{00}+T_{11})R+T_{01}=0.
\]
The TEP project outward-normal DtN is
\[
  \Lambda_{\mathrm{TEP}}=-(T_{00}+T_{01}R)=i\gamma.
\]

\section{Current Numerical Evidence}

The prototype \texttt{attempt/prototype/run\_all.m} checks the homogeneous
formula for \(M=8\), \(d=1\), \(L=0.73\), \(k_0=2.4\), and \(\beta=0.35\).
The first Octave sanity check produced project-DtN relative errors around
\(4\cdot10^{-17}\) and Riccati residuals around \(10^{-15}\).

\section{Limitations}

This draft does not yet contain:
\begin{itemize}
  \item obstacle-lead local cell DtN blocks;
  \item matrix Riccati or QZ implementation for non-diagonal blocks;
  \item Muller-DtN center coupling;
  \item comparison against the existing trace-subspace formulation;
  \item final bibliography verification for all sources.
\end{itemize}

\bibliographystyle{plain}
\bibliography{references}

\end{document}
